\section{Fluxo, Grafos Bipartidos e Reduções}

\subsection{Equivalências em Grafos Bipartidos}
Dado um grafo bipartido $G = (L \cup R, E)$ e o Pareamento Máximo $M$ (obtido via Fluxo Máximo $S \to L \to R \to T$ com capacidades $1$):

\begin{itemize}
    \item \textbf{Max Bipartite Matching (MBM):} Valor do fluxo máximo $|f_{\text{max}}|$.
    \item \textbf{Min Vertex Cover (MVC):} Mínimo de vértices para cobrir todas as arestas.
    \[ |\text{MVC}| = |\text{MBM}| \quad \text{(Teorema de König)} \]
    \item \textbf{Max Independent Set (MIS):} Máximo de vértices sem arestas entre si.
    \[ |\text{MIS}| = |V| - |\text{MVC}| = |V| - |\text{MBM}| \]
    \item \textbf{Min Path Cover (DAG):} Menor número de caminhos vértices-disjuntos para cobrir um DAG.
    Duplique cada vértice $u \to (u_L, u_R)$. Para cada aresta $u \to v$, adicione $(u_L, v_R)$ com capacidade $1$.
    \[ \text{Quantidade de Caminhos} = |V| - |\text{MBM}| \]
    \item \textbf{Teorema de Dilworth (Anticadeias em DAGs):} 
    Uma anticadeia é o maior conjunto de vértices onde nenhum alcança o outro.
    \[ \text{Tamanho da Maior Anticadeia} = \text{Quantidade Mínima de Caminhos} \]
\end{itemize}

\subsection{Teoremas em Grafos Gerais (Não Bipartidos)}
Enquanto o Teorema de König falha fora de grafos bipartidos, as \textbf{Identidades de Gallai} valem para \textbf{qualquer grafo} $G = (V, E)$ sem vértices isolados:

\begin{itemize}
    \item \textbf{Identidades de Gallai:}
    \[ |\text{MIS}| + |\text{MVC}| = |V| \]
    \[ |\text{Max Matching}| + |\text{Min Edge Cover}| = |V| \]

    \item \textbf{Min Edge Cover (Cobertura Mínima de Arestas):}
    Menor número de arestas necessárias para cobrir todos os vértices.
    \[ |\text{Min Edge Cover}| = |V| - |\text{Max Matching}| \]
    \textit{Como obter:} Encontre o Pareamento Máximo. Para cada vértice não pareado, adicione qualquer aresta incidente a ele.

    \item \textbf{Alertas para Programação Competitiva:}
    \begin{itemize}
        \item \textbf{König não vale:} $|\text{MVC}| \neq |\text{Max Matching}|$ em grafos gerais (ex: no triângulo $C_3$, $|\text{MVC}| = 2$ e $|\text{Matching}| = 1$).
        \item \textbf{Complexidade:} \text{MVC} e \text{MIS} tornam-se problemas \textbf{NP-Hard} em grafos gerais.
        \item \textbf{Fluxo não resolve Pareamento Geral:} Para pareamento em grafos arbitrários (com ciclos ímpares), o algoritmo de fluxo falha. É necessário usar o \textbf{Algoritmo de Blossom (Edmonds)}.
    \end{itemize}
\end{itemize}

\subsection{Como Recuperar a Resposta (Reconstrução)}

\textbf{Passo Inicial (Grafo Residual $G_f$):}
Rode uma DFS ou BFS a partir da fonte $S$ considerando apenas arestas com \textbf{capacidade residual $> 0$}. 

Guarde os vértices visitados no conjunto $S_{\text{reach}}$.

\begin{itemize}
    \item \textbf{Pareamento Máximo (MBM):}
    \begin{itemize}
        \item Uma aresta original $(u, v)$ com $u \in L$ e $v \in R$ pertence ao pareamento se tiver \textbf{fluxo positivo}.
    \end{itemize}

    \item \textbf{Min Cut (Arestas do Corte Mínimo):}
    \begin{itemize}
        \item O corte divide o grafo em visitados ($S_{\text{reach}}$) e não visitados ($T_{\text{reach}}$).
        \item Uma aresta $(u, v)$ do \textbf{grafo original} pertence ao corte se $u$ foi visitado e $v$ não foi visitado:
        \[ u \in S_{\text{reach}} \quad \text{e} \quad v \notin S_{\text{reach}} \]
        \item Essas arestas estarão obrigatoriamente \textbf{saturadas} ($\text{fluxo} = \text{capacidade}$).
    \end{itemize}

    \item \textbf{Min Vertex Cover (MVC em Bipartidos):}
    \begin{itemize}
        \item Pegue os vértices do lado esquerdo $L$ \textbf{NÃO visitados} pela DFS.
        \item Pegue os vértices do lado direito $R$ \textbf{VISITADOS} pela DFS.
        \[ \text{MVC} = (L \setminus S_{\text{reach}}) \cup (R \cap S_{\text{reach}}) \]
    \end{itemize}

    \item \textbf{Max Independent Set (MIS em Bipartidos):}
    \begin{itemize}
        \item É o \textbf{complemento} do Min Vertex Cover ($V \setminus \text{MVC}$):
        \[ \text{MIS} = (L \cap S_{\text{reach}}) \cup (R \setminus S_{\text{reach}}) \]
        \item Selecione $L$ \textbf{VISITADOS} e $R$ \textbf{NÃO visitados}.
    \end{itemize}

    \item \textbf{Min Path Cover (Recuperar Caminhos no DAG):}
    \begin{itemize}
        \item Reconstrua as cadeias seguindo as arestas $(u_L, v_R)$ que possuem $\text{fluxo} = 1$.
        \item Todo vértice $u_L$ que \textbf{não} tem fluxo saindo representa o \textbf{fim} de um caminho.
    \end{itemize}
\end{itemize}

\subsection{Reduções Clássicas de Fluxo}

\begin{itemize}
    \item \textbf{Project Selection (Maximum Weight Closure):}
    Queremos escolher um subconjunto de itens para maximizar o ganho líquido. Se um item $u$ for escolhido e depender de $v$ ($u \to v$), o item $v$ obrigatoriamente precisa ser escolhido junto.
    \begin{itemize}
        \item Item com \textbf{lucro} $p_u > 0$: aresta $S \to u$ com capacidade $p_u$.
        \item Item com \textbf{custo} $c_u > 0$: aresta $u \to T$ com capacidade $c_u$.
        \item \textbf{Dependência} $u \to v$: aresta $u \to v$ com capacidade $\infty$.
        \item \textbf{Corte Mínimo:} Representa a perda mínima.
        \item \textbf{Lucro Máximo:} $\sum (\text{Lucros Positivos}) - \text{Min Cut}$.
        \item \textbf{Itens Selecionados:} Conjunto $S_{\text{reach}} \setminus \{S\}$.
    \end{itemize}

    \item \textbf{Capacidade em Vértices:}
    Split no vértice $u \to (u_{\text{in}}, u_{\text{out}})$. Adicione aresta $u_{\text{in}} \to u_{\text{out}}$ com $\text{cap} = c(u)$. Arestas que entram em $u$ vão para $u_{\text{in}}$, e as que saem vão de $u_{\text{out}}$.
\end{itemize}

\subsection{Fluxo com Limite Inferior nas Arestas (Circulation with Demands)}
Use essa técnica quando uma aresta $u \to v$ possuir uma capacidade mínima obrigatória $L$ e uma máxima $R$ ($L \le f \le R$). O algoritmo tradicional não entende limites inferiores, então precisamos usar vértices fictícios para forçar esse fluxo.

\textbf{Construção da Rede:}
\begin{itemize}
    \item Crie dois vértices auxiliares: uma \textbf{Super Fonte ($SS$)} e um \textbf{Super Sumidouro ($TT$)}.
    \item Para cada aresta original $u \to v$ com limite $[L, R]$, adicione na nova rede uma aresta $u \to v$ com \textbf{capacidade $R - L$}.
    \item Mantenha um balanço para cada vértice $v$. Seja $in[v]$ a soma de todos os limites $L$ das arestas que \textbf{entram} em $v$, e $out[v]$ a soma de todos os limites $L$ das arestas que \textbf{saem} de $v$.
    \item Adicione uma aresta $SS \to v$ com capacidade $in[v]$.
    \item Adicione uma aresta $v \to TT$ com capacidade $out[v]$.
    \item Adicione uma aresta do sumidouro original para a fonte original ($T \to S$) com capacidade $\infty$ (para permitir a circulação).
\end{itemize}

\textbf{Verificação e Resposta:}
\begin{itemize}
    \item Rode o Fluxo Máximo partindo de $SS$ para $TT$.
    \item \textbf{Existe solução válida?} Apenas se \textbf{todas} as arestas saindo de $SS$ (ou chegando em $TT$) estiverem totalmente saturadas. Ou seja, se o $\text{Max Flow} = \sum in[v]$.
    \item \textbf{Recuperar fluxo nas arestas:} O fluxo real em uma aresta original $u \to v$ será: \\
    $L + (\text{fluxo que passou pela aresta } u \to v \text{ na nova rede})$.
    \item \textbf{Para achar o Max Flow real de $S$ a $T$:} Após verificar que o fluxo inicial é válido, remova os vértices $SS$ e $TT$ e suas arestas, remova a aresta infinita $T \to S$, e rode mais uma vez o Fluxo Máximo de $S$ para $T$ na rede residual.
\end{itemize}
